\section{Discussion}

The population-size threshold is not the surprising part of this result --
that bigger populations help on harder landscapes is close to the null
expectation. The surprising part is that the specific, quantitative,
textbook mechanism for \emph{why} -- shrinking per-mutation selection
strength letting drift dominate, exactly as Kimura's formula describes for
a single locus in a real population -- points the wrong way. NK-model
epistasis does not appear to dilute the effect of a single beneficial
mutation; if anything, in landscapes with more interacting loci, a bit
flip that does help tends to help by more, not less, on average.

A more plausible mechanism, consistent with what is already known about
NK landscapes -- that the number of local optima grows with $K$
\cite{kauffman1987towards,weinberger1991local} -- is that the binding
constraint at high $K$ is not single-locus fixation probability but
population-level exploration: a small population converges onto one local
optimum's basin before crossover and mutation can discover a better one
elsewhere, while a larger population maintains enough simultaneous
diversity to explore multiple basins in parallel. This is a
population-\emph{diversity} argument, not a population-\emph{genetics}
one in the classical single-locus sense, and it was not tested directly
here (it would require tracking, e.g., how many distinct basins of
attraction a population's genomes span over time). We flag it as the
natural next hypothesis rather than claim it is confirmed.

A second limitation is resolution: the local pilot's 4-point $K$ grid gave
a genuinely weak test of linear-vs-log-linear scaling. The GPU-scale run's
10-point grid is a materially stronger test of the same question, and it
reaches the same conclusion (linear preferred by AIC) -- so this is not
merely "the pilot was underpowered and the real answer is log-linear
after all," in the spirit of this repository's project 06 pilot, but a
result that held up under more scrutiny, not one that dissolved under it.
The non-monotonicity in the fitted $N_c(K)$ at high $K$, however, did not
go away with more landscape averaging (10 draws per $K$ at GPU scale vs 5
locally), which suggests either landscape-instance variance at high $K$
needs substantially more than 10 draws to average out, or the true
$N_c(K)$ relationship itself is not smooth -- we cannot distinguish these
with the current data and report the non-monotonicity as-is rather than
picking whichever explanation is more convenient.
